\documentclass[11pt]{article}
\usepackage[margin=1in]{geometry}
\usepackage{hyperref}
\begin{document}
\begin{center}\large Qu et al. (2022): full-text acquisition pending\end{center}
\noindent\textbf{[Unread].} This is an acquisition record, not a scientific summary.
The full paper and supplement have been requested from the user.

\paragraph{Reference.} Qu et al., ``Attribution of the 2020 surge in atmospheric
methane by inverse analysis of GOSAT observations,'' \emph{Environmental Research
Letters} 17, 094003,
\href{https://doi.org/10.1088/1748-9326/ac8754}{doi:10.1088/1748-9326/ac8754}.
The original bibliography is retained unchanged in the repository.

\paragraph{Acquisition attempts.} The IOP article and PDF endpoints returned
access challenges. The IOP content host returned HTTP 403 for the article and
supplementary-data URLs. OpenAlex and author-publication pages were checked
for an accessible copy. A NASA Earth Science Project Office publication entry
was retrieved, but its current page supplies a DOI link without an attached
PDF. Download attempts are recorded in \texttt{data/downloads.jsonl}.

\paragraph{Required material.} The complete article and all supplementary files
are required before documenting its inversion design, posterior fluxes, OH
definition, uncertainty ensemble, or regional attribution. A summary of this
study in a later paper is not a substitute for reading its primary methods.
No values attributed to this article are inserted into the numerical comparison.

\paragraph{Planned use once supplied.} Extract native 2019 and 2020 emissions,
OH loss frequencies or matched burden and removal fluxes, soil and other loss
terms, and the averaging conventions used to report growth. Rebase all
available quantities to 2019 and apply the documented He et al. equation.
Preserve differences between annual priors, initialization, and observation
sampling when comparing it with Qu et al. (2024). Do not treat the latter as
an interchangeable replacement for this unavailable study, and do not infer
missing 2021--2022 values by extrapolation.
\end{document}
